\section*{Mission : Interception des pirates }
\setcounter{chapter}{1}

Notre équipage a détecté une anomalie dans les communications de la flotte pirate. 
Les données interceptées sont fragmentées : nous avons le code des fonctions, mais leurs noms ont été effacés pour des raisons de sécurité.
Votre mission : analyser ces signaux et restaurer les protocoles de communication pour notre équipage.

\begin{Definition}[Le Schéma d'une Fonction ]
Une fonction est un protocole de traitement de données. Elle reçoit des \motcle{paramètres} (entrées), utilise des \motcle{variables locales} pour ses calculs internes, et renvoie un \motcle{retour} (résultat final).
\end{Definition}

\section{Audit de signaux impériaux : Comprendre }

\begin{Exercise}[title={Signal Alpha (Simple) }]
\begin{lstlisting}[style=pythonstyle, language=python]
def fonction_mystere(vitesse_hyperespace):
    temps_voyage = 5000 / vitesse_hyperespace
    return temps_voyage

calcul = fonction_mystere(50)
\end{lstlisting}
\begin{enumerate}
    \item Quel est le paramètre (entrée) et son type ?
    \item Quelle est la variable locale et son type ?
    \item Quel est le retour (sortie) et son type ?
    \item Entourez la définition de la fonction et soulignez l'appel de la fonction.
    \item Si la flotte pirate envoie une \texttt{vitesse\_hyperespace} de 100, quelle valeur numérique exacte sera stockée dans la variable calcul ?
    \item \textbf{Question :} Quel nom donneriez-vous à cette fonction ? 
\end{enumerate}
\end{Exercise}

\begin{Exercise}[title={Signal Beta (Simple) }]
\begin{lstlisting}[style=pythonstyle, language=python]
def fonction_mystere(niveau_alerte):
    if niveau_alerte > 1000:
        est_critique = True
    else:
        est_critique = False
    return est_critique

test = fonction_mystere(1200)
\end{lstlisting}
\begin{enumerate}
    \item Quel est le paramètre (entrée) et son type ?
    \item Quelle est la variable locale et son type ?
    \item Quel est le retour (sortie) et son type ?
    \item Entourez la définition de la fonction et soulignez l'appel de la fonction.
    \item Que retourne la fonction pour \texttt{niveau\_alerte = 100} et \texttt{niveau\_alerte = 5000} ?
    \item \textbf{Question :} Quel nom donneriez-vous à cette fonction ? 
\end{enumerate}
\end{Exercise}

\begin{Exercise}[title={Cryptage de la station (Complexe) }]
\begin{lstlisting}[style=pythonstyle, language=python]
def fonction_mystere(code_acces):
    code_inverse = ""
    for caractere in code_acces:
        code_inverse = caractere + code_inverse
    return code_inverse

decryptage = fonction_mystere("AB421")
\end{lstlisting}
\begin{enumerate}
    \item Quel est le paramètre (entrée) et son type ?
    \item Quelles sont les variables locales et leurs types ?
    \item Quel est le retour (sortie) et son type ?
    \item Entourez la définition de la fonction et soulignez l'appel de la fonction.
    \item Que retourne la fonction pour \texttt{code\_acces = "PIRATE"} ?
    \item Si l'entrée est 'PIRATE', quel est le tout premier caractère placé dans la variable \texttt{code\_inverse} ? Quel est le dernier ?
    \item \textbf{Question :} Quel nom donneriez-vous à cette fonction ? 
\end{enumerate}
\end{Exercise}

\begin{Exercise}[title={Scanner de robots (Complexe) }]
\begin{lstlisting}[style=pythonstyle, language=python]
def fonction_mystere(inventaire_pieces):
    pieces_valides = []
    for piece in inventaire_pieces:
        if "cyber" in piece:
            pieces_valides.append(piece)
    return pieces_valides

resultat = fonction_mystere(["cyber-bras", "boulon", "cyber-oeil"])
\end{lstlisting}
\begin{enumerate}
    \item Quel est le paramètre (entrée) et son type ?
    \item Quelles sont les variables locales et leurs types ?
    \item Quel est le retour (sortie) et son type ?
    \item Entourez la définition de la fonction et soulignez l'appel de la fonction.
    \item Le signal intercepté est ["moteur", "cyber-jambe", "laser"]. Quelle liste la fonction va-t-elle renvoyer ?
    \item \textbf{Question :} Quel nom donneriez-vous à cette fonction ? 
\end{enumerate}
\end{Exercise}

\newpage
\section{Protocoles de l'équipage : Rédiger }

\begin{Exercise}[title={La fonction \texttt{calculer\_carburant} (Simple) }]
\textbf{Description :} L'équipage doit prévoir le ravitaillement de ses vaisseaux. Chaque moteur consomme exactement 15 litres de carburant pour parcourir un parsec. Votre mission est de créer un programme qui, à partir d'une distance donnée, calcule le volume total de carburant nécessaire pour la mission.
\begin{enumerate}
    \item Explicitez l'entrée et la sortie.
    \item \textbf{Exemple :} Si l'entrée est \texttt{10}, la sortie est \dotfill.
    \item Donner une entrée pour la sortie 30000.
    \item \textbf{Logique :} Complétez :
    \begin{itemize}
        \item Je commence par définir la fonction nommée \dotfill
        \item Cette fonction accepte un paramètre d'entrée appelé \dotfill
        \item À l'intérieur de la fonction, je crée une variable locale appelée \dotfill
        \item Pour calculer sa valeur, je multiplie la distance reçue en entrée par le coefficient de consommation qui est de \dotfill
        \item Enfin, j'utilise l'instruction \dotfill pour renvoyer la valeur contenue dans la \dotfill
    \end{itemize}
    
\end{enumerate}
\end{Exercise}

\begin{Exercise}[title={La fonction \texttt{est\_allie} (Simple) }]
\textbf{Description :} Pour éviter les tirs amis, les systèmes de défense doivent identifier les matricules des vaisseaux. On considère qu'un vaisseau appartient à notre flotte si son matricule commence par les caractères "RX" ou "CZ". La fonction doit répondre par "Vrai" ou "Faux".
\begin{enumerate}
    \item Explicitez l'entrée (chaîne \texttt{matricule}) et la sortie (booléen).
    \item \textbf{Exemple :} \\- Si l'entrée est \texttt{RXD6}, la sortie est \dotfill.\\
    - Si l'entrée est \texttt{D6F4}, la sortie est \dotfill.
    \item Donner une entrée pour la sortie True.
    \item Donner une entrée pour la sortie False.
    \item \textbf{Logique :} Complétez :
    \begin{itemize}
        \item Je définis la fonction \dotfill avec pour paramètre la variable \dotfill.
        \item J'utilise une structure de condition \dotfill pour vérifier le début de la chaîne.
        \item Le test consiste à vérifier si les deux premiers caractères du matricule sont égaux à \dotfill ou s'ils sont égaux à \dotfill.
        \item Si le test est positif, je choisis de renvoyer immédiatement la valeur \dotfill
        \item Dans le cas contraire (sinon), je renvoie la valeur \dotfill.
    \end{itemize}
    \item \textbf{À vous de coder :} Rédigez la fonction complète.
\end{enumerate}
\end{Exercise}

\newpage
\begin{Exercise}[title={La fonction \texttt{sceller\_donnees} (Complexe) }]
\textbf{Description :} Les pirates tentent d'écouter nos fréquences radio. Pour les protéger, nous allons appliquer un "masque" numérique. La fonction prend une liste de fréquences (nombres) et une clé secrète. Elle doit transformer chaque fréquence de la liste en lui ajoutant la valeur de la clé.
\begin{enumerate}
    \item Identifiez les entrées (liste de fréquences et clé entière) et la sortie.
    \item \textbf{Exemple :} Pour \texttt{[100, 200]} et une clé de \texttt{7}, la sortie est \texttt{[107, 207]}.
    \item Donner les entrées pour la sortie \texttt{[23, 33, 43, 63]} (plusieurs réponses possibles).
    \item Donner les entrées pour la sortie \texttt{[7, 14, 9]} (plusieurs réponses possibles).
    \item \textbf{Logique :} Complétez :
    \begin{itemize}
        \item Je définis la fonction \dotfill avec ses deux entrées : \dotfill et \dotfill.
        \item Au début, je crée une variable locale nommée \dotfill qui contient une liste \dotfill.
        \item J'installe un parcours par \dotfill afin de parcourir un par un tous les éléments de la liste \dotfill.
        \item À l'intérieur de la boucle, je crée une variable temporaire qui est égale à la valeur de la fréquence actuelle \dotfill à la \dotfill.
        \item J'utilise ensuite une méthode pour \dotfill ce nouveau nombre à la fin de ma liste \dotfill.
        \item Une fois que la boucle a fini de traiter tous les éléments, je renvoie la liste \dotfill complète.
    \end{itemize}
    \item \textbf{À vous de coder :} Rédigez la fonction complète.
\end{enumerate}
\end{Exercise}

\begin{Exercise}[title={La fonction \texttt{code\_cesar\_un} (Complexe) }]
\textbf{Description :} Pour crypter un texte simple, nous utilisons le décalage d'un rang : chaque lettre est remplacée par la lettre suivante dans l'alphabet (le 'a' devient 'b', le 'b' devient 'c', etc.).
\begin{enumerate}
    \item Entrée : \texttt{mot} (String). Sortie : \texttt{mot\_code} (String).
    \item \textbf{Exemple :} \\- Si l'entrée est \texttt{Bonjour tout le monde}, la sortie est \dotfill.\\
    - Si l'entrée est \texttt{Je suis ton ami}, la sortie est \dotfill.
    \item Le message intercepté est \texttt{mb dibodf ftu fo upj}. Sachant que chaque lettre a été décalée d'un rang vers la droite (a devient b), retrouvez le message original.
    \item Donner les entrées pour la sortie \texttt{qfuju b qfuju m pjtfbv gbju tpo oje}.
    \item \textbf{Logique :} Complétez :
    \begin{itemize}
        \item Je définis la fonction \dotfill avec le paramètre \dotfill.
        \item Je déclare une variable locale nommée \dotfill qui contient toutes les lettres de "a" à "z".
        \item J'initialise une variable \dotfill avec \dotfill.
        \item Je lance une boucle pour examiner chaque \dotfill présente dans le \dotfill d'origine.
        \item Pour chaque lettre, je cherche d'abord son \dotfill (sa position numérique) dans la variable \dotfill.
        \item Je calcule la position suivante en ajoutant \dotfill à cet indice.
        \item Je récupère la lettre correspondante à ce nouvel indice dans l'alphabet et je l'\dotfill (concaténation) à ma chaîne \dotfill.
        \item Je termine en renvoyant le \dotfill final.
    \end{itemize}
    \item \textbf{À vous de coder :} Rédigez la fonction complète.
\end{enumerate}
\end{Exercise}

\newpage

\section*{Exercices Bonus : Protocoles de Haute Sécurité}

\begin{Exercise}[title={Bonus 1 : Le Grand Décalage de César (Complexe) }]
\textbf{Description :} C'est une version améliorée du code précédent. Au lieu de décaler toujours de 1, l'utilisateur peut choisir la valeur du décalage (n). Si on dépasse la lettre 'z', on doit recommencer au début de l'alphabet (modulo 26).
\begin{enumerate}
\item \textbf{Entrées :} Identifiez les deux entrées nécessaires (le mot et la valeur du décalage).
\item \textbf{Exemple :} Le signal de sortie est 'sur'. Sachant que la flotte pirate a utilisé un décalage n = 2, quel était le message d'origine avant le cryptage ?
\item \textbf{Rétro-ingénierie :} Si la sortie est \texttt{"sur"} avec un décalage de \texttt{2}, quel était le mot original ?
\item \textbf{Logique :} Complétez :
    \begin{itemize}
        \item Je définis la fonction avec deux paramètres : \dotfill et \dotfill.
        \item Comme précédemment, je prépare une variable pour le résultat et je possède la chaîne de l'alphabet.
        \item Je parcours chaque lettre du mot avec une \dotfill.
        \item Pour chaque lettre, je trouve son \dotfill, puis je lui ajoute la valeur \dotfill.
        \item Pour éviter de sortir de l'alphabet (si on dépasse 25), j'utilise l'opération mathématique \dotfill sur le nouvel \dotfill.
        \item Je cherche la \dotfill située à ce nouvel indice "sécurisé" et je l'\dotfill dans ma variable de sortie.
        \item Je renvoie le résultat une fois la boucle terminée.
    \end{itemize}
\item \textbf{À vous de coder :} Rédigez la fonction complète nommée \texttt{cesar\_general}.
\end{enumerate}
\end{Exercise}

\begin{Exercise}[title={Bonus 2 : Le Chiffre de Vigenère (Expert) }]
\textbf{Description :} C'est le niveau ultime. On n'utilise plus un chiffre pour décaler, mais une liste de décalages qui tourne en boucle (ex: \texttt{[1, 2]} pour décaler la 1ère lettre de 1, la 2ème de 2, la 3ème de 1, etc.). 
\begin{enumerate}
\item \textbf{Entrées :} Identifiez les deux entrées (le mot et la liste des décalages).
\item \textbf{Exemple :} Pour le mot \texttt{"aaaa"} et la liste \texttt{[1, 2, 3]}, la sortie est \texttt{"bcbd"}.
\item \textbf{Logique :} Complétez :
    \begin{itemize}
        \item Je définis la fonction avec ses deux paramètres \dotfill et \dotfill.
        \item Je crée une variable locale appelée \dotfill (un compteur) initialisée à \dotfill pour suivre ma position dans la liste des décalages.
        \item Je parcours le mot lettre par lettre grâce à une \dotfill.
        \item À chaque étape, je récupère le décalage situé à l'indice \dotfill dans la \dotfill.
        \item J'applique ce décalage à la lettre actuelle (en utilisant le modulo 26 pour la sécurité).
        \item Après avoir ajouté la lettre codée au résultat, j'augmente mon compteur \dotfill de \dotfill.
        \item J'ajoute une condition : si mon compteur i devient égal à la \dotfill, alors je remets \dotfill à \dotfill pour recommencer au début de la clé.
        \item Je renvoie le message crypté \dotfill final.
    \end{itemize}
\item \textbf{À vous de coder :} Rédigez la fonction complète nommée \texttt{vigenere\_liste}.
\end{enumerate}
\end{Exercise}